\documentclass[12pt,a4paper]{ctexbook}
\usepackage{geometry}
\geometry{margin=2.2cm}
\usepackage{setspace}
\setstretch{1.35}
\usepackage{hyperref}
\title{全宋诗选·poet.song.3000}
\author{古典文献整理}
\date{}
\begin{document}
\maketitle
\tableofcontents
\newpage
\section{五哀詩 故殿中侍御史滎陽鄭公}
\textbf{王禹偁}\par\vspace{0.5em}
柱史有名跡，清才自天縱。\par
構思慶雲合，落筆醴泉湧。\par
歌詩與文賦，錚錚人口諷。\par
揚袂入澤宮，鵠心一箭中。\par
恃才善戲謔，負氣好侮弄。\par
大志有誰知，細行乖自訟。\par
小諫事世宗，惕惕佩光寵。\par
太祖方歷試，握兵權已重。\par
上書范魯公，先見不能用。\par
曆數不在周，謳謠卒歸宋。\par
汗漫失屠龍，接輿遂歌鳳。\par
行荷伯倫鍤，高卧畢卓甕。\par
神德不爲嫌，優待臺諫俸。\par
晚求萬泉令，吏隠官資冗。\par
一旦隨朝露，識者彌哀痛。\par
無子嗣家聲，身世若一夢。\par
文編多散失，人口時傳誦。\par
空持一器酒，何處澆孤冢。\par

\section{五哀詩 故國子博士郭公}
\textbf{王禹偁}\par\vspace{0.5em}
汾陽飽經術，賦性甚坦率。\par
在昔舉神童，廣場推傑出。\par
尚書誦在口，何論落自筆。\par
總角取科名，弱冠紆纓紱。\par
早佐湘陰幕，漢鼎入周室。\par
失志罷屠龍，佯狂遂捫蝨。\par
周行亦黽勉，吏隠多放逸。\par
滑稽東方朔，圖畫王摩詰。\par
古文識科斗，奧學辨萍實。\par
字窮蒼頡本，篆證陽冰失。\par
王績醉爲鄉，伯倫居無匹。\par
俸錢乏一囊，宦路從三黜。\par
朱衣多不著，白髮仍慵櫛。\par
漸老覊旅年，方見昇平日。\par
忽以伎術召，此意殊鬰鬰。\par
放口忤無鬚，何門求造膝。\par
遁逃終見捕，譴逐道中卒。\par
遺孤落閭閻，荒冢鳴蟋蟀。\par
手澤漸難求，誰家耀箱帙。\par
投弔焚此詩，九原應有物。\par

\section{五哀詩 故太子中允知洛陽縣潁公}
\textbf{王禹偁}\par\vspace{0.5em}
洛陽富文彩，峭拔四子流。\par
提筆入廣場，辭氣干斗牛。\par
擢第在芸閣，言事觸冕旒。\par
左降宰百里，道勝心無憂。\par
才高恥吏役，放蕩不檢脩。\par
起應賢良科，下筆不見休。\par
青宮尚淹恤，赤縣且優游。\par
輕財糞土賤，高義雲天浮。\par
懸磬任貧窶，盈樽長獻酬。\par
知己彼何人，鳳閣與鼇頭。\par
推挽終不起，壯志將焉收。\par
晚年圯橋役，關市良可羞。\par
忽然爲異物，寒草封一丘。\par
孀嫠應凍餓，交友誰尋求。\par
遺孫方稚齒，爽秀已遒遒。\par
皇天若有憑，必使光貽謀。\par

\section{太一宮祭迴馬上偶作寄韓德純道士}
\textbf{王禹偁}\par\vspace{0.5em}
去年暨今夏，承詔祠天神。\par
昔當摇落時，宮葉紅紛紛。\par
此來芳春暮，宮草青蓁蓁。\par
三日奉齋潔，百骸祛垢氛。\par
御署致恭虔，天香何絪縕。\par
監祀黄門郎，攝官紫微臣。\par
松枝拂劍佩，樹影光搢紳。\par
燈燭晃白晝，香花藹清芬。\par
金殿禮神仙，瑶壇醮星辰。\par
至誠不爲己，玄鑒當福仁。\par
質明祀事畢，復命趨紫宸。\par
揚鞭入村落，緩轡聊逡巡。\par
麥田少時雨，蠶月無閒人。\par
自慚懷祿仕，蠹此力穡民。\par
又拋三洞趣，來入九衢塵。\par
不如韓道士，長此養天真。\par

\section{送筇杖與劉湛然道士}
\textbf{王禹偁}\par\vspace{0.5em}
有客遺竹杖，九節共一枝。\par
鶴脰老更長，龍骨乾且奇。\par
我問何所來，來從西南夷。\par
因思漢武帝，求此民力疲。\par
明明聖天子，德教加四維。\par
蠻貊盡臣妾，縣道皆覊縻。\par
僰僮與笮馬，入貢何纍纍。\par
此竹日已賤，輕視如蒿藜。\par
我年三十七，血氣未全衰。\par
况在紫微垣，動爲簪笏覊。\par
倚壁如長物，歲月無所施。\par
寸心空愛惜，惜此來天涯。\par
忽承明主詔，來謁太一祠。\par
再見劉先生，氣貌清且羸。\par
持此以爲贈，所謂得其宜。\par
少助橘童力，好引花鹿隨。\par
步月莫離手，看山聊搘頤。\par
微物懶致書，故作筇竹詩。\par

\section{對雪}
\textbf{王禹偁}\par\vspace{0.5em}
帝鄉歲云暮，衡門晝長閉。\par
五日免常參，三館無公事。\par
讀書夜卧遲，多成日高睡。\par
睡起毛骨寒，窗牖瓊花墜。\par
披衣出戶看，飄飄滿天地。\par
豈敢患貧居，聊將賀豐歲。\par
月俸雖無餘，晨炊且相繼。\par
薪芻未闕供，酒肴亦能備。\par
數杯奉親老，一酌均兄弟。\par
妻子不饑寒，相聚歌時瑞。\par
因思河朔民，輸稅供邊鄙。\par
車重數十斛，路遥幾百里。\par
羸蹄凍不行，死轍冰難曳。\par
夜來何處宿，闃寂荒陂裏。\par
又思邊塞兵，荷戈禦胡騎。\par
城上卓旌旗，樓中望烽燧。\par
弓勁添氣力，甲寒侵骨髓。\par
今日何處行，牢落窮沙際。\par
自念亦何人，偷安得如是。\par
深爲蒼生蠹，仍尸諫官位。\par
謇諤無一言，豈得爲直士。\par
褒貶無一詞，豈得爲良史。\par
不耕一畝田，不持一隻矢。\par
多慚富人術，且乏安邊議。\par
空作對雪吟，勤勤謝知己。\par

\section{金吾}
\textbf{王禹偁}\par\vspace{0.5em}
金吾河朔人，事郡在賤列。\par
攀附周世宗，龍飛起魚鱉。\par
委質向聖朝，積功取旄鉞。\par
所在肆貪殘，乘時恃勳伐。\par
皇家平金陵，九江聚遺孽。\par
彌年城乃陷，不使鷄犬活。\par
老小數千人，一怒盡流血。\par
三惑無不具，五福何嘗缺。\par
晚年得執金，富貴居朝闕。\par
娛樂有清商，康强無白髮。\par
享年六十九，固不爲夭折。\par
考終北牖下，手足全啟發。\par
子孫十數人，解珮就衰絰。\par
贈典頗優崇，視朝爲之輟。\par
哀榮既如是，報應何足說。\par
責簿李廣死，賜劍武安滅。\par
僥倖過古人，况無大功烈。\par
福善與禍淫，斯言僅虛設。\par
吁嗟爲儒者，寒餓守名節。\par
五十朝大夫，龍鍾頭似雪。\par
無故不殺羊，禮文安可越。\par
何况賓祀間，貧苦無羊殺。\par

\section{送戚維戚綸之閬州亳州}
\textbf{王禹偁}\par\vspace{0.5em}
古人貴道德，豈以祿位拘。\par
有道不在位，顔回舜之徒。\par
無德殃且至，商受爲獨夫。\par
以此固名節，富貴安足圖。\par
睢陽戚先生，今世古之儒。\par
終身不求仕，沒齒唯誦書。\par
孝愛睦姻族，淳謹化里閭。\par
文翁變蜀土，尚以德位俱。\par
仲淹居汾水，專用素教敷。\par
青青子衿輩，櫛比曳朝裾。\par
盛德苟無報，吾道將焉如。\par
鯉庭生二子，驪頷委雙珠。\par
樸厚有父風，學業張皇謨。\par
各登進士科，齊乘使者車。\par
閬中佐郡政，亳社司搉酤。\par
佐郡古半刺，可使惸嫠蘇。\par
搉酤今權道，亦助軍國須。\par
牽車驥雖屈，補履刃有餘。\par
誰謂佐著作，在上積薪芻。\par
誰謂光祿勳，其中混瓊琚。\par
我亦何爲者，先上青雲衢。\par
伐木空求友，拔茅未連茹。\par
終待奮直筆，會當伏青蒲。\par
請慎名與器，願分賢與愚。\par
彙征補皇極，戮力張道樞。\par
巨魚方喣沫，相忘在江湖。\par
冥鴻當順風，接翼摩太虛。\par
此志尚壹鬰，此別良悲吁。\par
苦縣且密邇，劍關多崎嶇。\par
于役勿爲歎，加飯聊自娛。\par
尺素惠好音，無惜雙鯉魚。\par

\section{送陳侯之任同州}
\textbf{王禹偁}\par\vspace{0.5em}
丈夫貴自奮，何必恃貽謀。\par
當年功不立，古人重爲憂。\par
太師忠順公，沒世有餘休。\par
令子一何賢，恥隨資蔭流。\par
碌碌奉朝請，帝恩何以酬。\par
鬰鬰不快意，上章聞冕旒。\par
願爲循良吏，未甘恩澤侯。\par
天子嘉乃誠，錫命蒞同州。\par
同州古三輔，北望雄且優。\par
詔條得以布，民瘼得以求。\par
勿謂綺紈子，當有襦袴謳。\par
詩云無忝爾，禮云必爲裘。\par
去去繼旄節，免作將門羞。\par

\section{送朱九齡}
\textbf{王禹偁}\par\vspace{0.5em}
吏隠不求貴，親老不擇祿。\par
之子有俊才，弱冠中正鵠。\par
弗墜先人業，何慚有道穀。\par
一命佐碭山，枳棘聊容足。\par
再命宰嘉興，絲桐幾易俗。\par
率身甘菜茹，養母求粱肉。\par
承顔苟不虧，折腰未爲辱。\par
解印無餘貲，舟中只琴筑。\par
十口寄淮泗，一身來輦轂。\par
又說東南行，秋風江水淥。\par
鄱陽古名郡，赤金流山谷。\par
每歲鼓錢刀，從來設官局。\par
還得便高堂，無辭縻逸躅。\par
颺帆江湖思，木脫天地肅。\par
楓葉紫斕斒，蓼花紅\{罒/彔\}𦌊。\par
津吏曉來迎，溪僧夜留宿。\par
至止事方簡，優游從所欲。\par
江城豐稻粱，水市多魚蔌。\par
三載奉甘飴，百錢飽家族。\par
自有綵衣華，勿歎藍袍綠。\par
行年未三十，氣壯顔如玉。\par
行義日以聞，焉能長碌碌。\par
終列侍臣班，耀我同年錄。\par
且賦白華詩，唱作離筵曲。\par

\section{暴富送孫何入史館}
\textbf{王禹偁}\par\vspace{0.5em}
孟郊常貧苦，忽吟不貧句。\par
爲喜玉川子，書船歸洛浦。\par
迺知君子心，所樂在稽古。\par
漢公得高科，不足唯墳素。\par
二年佐棠陰，眼黑怕文簿。\par
躍身入三館，爛目閱四庫。\par
孟貧昔不貧，孫貧今暴富。\par
暴富亦須防，文高被人妬。\par

\section{贈劉仲堪}
\textbf{王禹偁}\par\vspace{0.5em}
劉生頗少秀，爲學識根柢。\par
丘軻有堂奧，試脚到階砌。\par
楊墨恣荒榛，揮手欲芟薙。\par
携文訪謫居，趣向非權勢。\par
對挹雛鳳下，交言孤鶴唳。\par
在璞認良玉，行當爲國器。\par
彼茁見靈芽，佇可供王祭。\par
豈止隨衆人，區區一枝桂。\par
宜哉孫漢公，妻之以女弟。\par
吾家兄之子，笄年未伉儷。\par
恨不早相逢，取子爲佳壻。\par

\section{送馮尊師}
\textbf{王禹偁}\par\vspace{0.5em}
前日訪潘閬，下馬入窮巷。\par
忽見雙笋石，卧向青苔上。\par
云是馮尊師，秋來留在茲。\par
今說東南行，問我堅乞詩。\par
又見宋閣老，亦言師甚好。\par
欲去天台山，即別長安道。\par
臺閣有羣英，贈別瑰與瓊。\par
琤然滿懷袖，此事殊爲榮。\par
安用徵吾句，吾道方齟齬。\par
老爲八品官，有山未能去。\par
束髮號男兒，出處貴得宜。\par
出則學臯夔，獨立稱帝師。\par
處則同喬松，决起如冥鴻。\par
誰能似蚯蚓，蟠屈泥土中。\par
師行甚可羨，雲鶴無覊絆。\par
爲我持此詩，題于桐栢觀。\par

\section{八絕詩 庶子泉}
\textbf{王禹偁}\par\vspace{0.5em}
物形固天造，物景不自勝。\par
泉乎未遇人，石罅徒流迸。\par
宮相政多暇，行樂躡巖磴。\par
發蒙漲爲溪，幽致茲焉盛。\par
唐賢大曆後，峭壁刻名姓。\par
我來一何暮，今秋始乘興。\par
山勢環有缺，山門壺引柄。\par
乍挹清泚姿，頗愜幽閑性。\par
味將春茗宜，光與曉嵐暝。\par
架竹落僧厨，遠聲入清磬。\par
何當宿禪室，欹枕終夜聽。\par
飲多病骨換，照久塵逈。\par
銷盡謫居愁，無心治歸艇。\par

\section{八絕詩 白龍泉}
\textbf{王禹偁}\par\vspace{0.5em}
崖石何鑿鑿，渤潏湧山脚。\par
含虛光可鑒，倒影壁如削。\par
瀰漫到前溪，支脈通遠壑。\par
白龍乃神物，胡此戀一勺。\par
勢能爲霖雨，力可興雷雹。\par
上天未有命，茲焉縮頭角。\par
應防困螻蟻，泥蟠詎敢躍。\par
葅醢曾是憂，豢養殊非樂。\par
亦如君子道，出處貴先覺。\par
吾生苦遷謫，未免郡政縛。\par
偶來泉畔坐，感興題秋籜。\par
斗藪纓上塵，試就清漣濯。\par
聊將一掬水，洗我面慚怍。\par
日暮不忍歸，風埃滿城郭。\par

\section{八絕詩 明月溪}
\textbf{王禹偁}\par\vspace{0.5em}
漲溪者爲誰，人骨皆已朽。\par
我來尋故跡，溪荒亂泉吼。\par
惜哉幽勝事，盡落唐賢手。\par
唯餘舊時月，團團照山口。\par

\section{八絕詩 清風亭}
\textbf{王禹偁}\par\vspace{0.5em}
茲亭廢已久，厥趾猶在哉。\par
清風爲我起，疑有精靈來。\par
神交念宮相，臨砌傾一杯。\par
迴頭問黄菊，寂寞與誰開。\par

\section{八絕詩 望日臺}
\textbf{王禹偁}\par\vspace{0.5em}
荒臺隠層碧，雲磴踰百尺。\par
攀蘿試一上，依然有遺跡。\par
掌舒舊砌平，屏卓諸峰直。\par
憑高聊寫望，孤懷念鄉國。\par
長安不可見，但對金烏赤。\par
傾盡葵藿心，庶免浮雲隔。\par

\section{八絕詩 歸雲洞}
\textbf{王禹偁}\par\vspace{0.5em}
碧洞何眈眈，呀然倚山根。\par
朝雲出如嘔，暮雲歸如吞。\par
大塊氣不死，茲爲玄牝門。\par
怪石擁左右，勢若貙虎蹲。\par
旁行數十步，漆黑不可捫。\par
安得鞭燭龍，爲吾前馳奔。\par
尋盡神仙窟，一一知其源。\par

\section{八絕詩 陽冰篆}
\textbf{王禹偁}\par\vspace{0.5em}
泠泠庶子泉，落落陽冰筆。\par
雲氣勢崩垂，龍蛇互蟠屈。\par
嶧山既劘滅，石鼓又缺失。\par
唯茲數十字，遒勁倚雲窟。\par
模印徧華夷，流傳耀緗帙。\par
書誠一藝爾，小道詎可忽。\par
乃知出人事，千古名不沒。\par

\section{八絕詩 垂藤蓋}
\textbf{王禹偁}\par\vspace{0.5em}
古藤何樛蟠，低蔭庶子泉。\par
童童若青蓋，挂在絕壁前。\par
月穿波影碎，露滴岸聲圓。\par
曉籠使君醉，夜覆溪僧禪。\par
昔之好事者，爲作出偶然。\par
人亡藤已朽，猶得聲名傳。\par
我欲移青松，植向泉眼邊。\par
既圖歲月久，復取節操全。\par
所爲亦一時，所期或千年。\par
庶吾不朽名，長與偃蓋懸。\par

\section{東門送郎吏行寄承旨宋侍郎}
\textbf{王禹偁}\par\vspace{0.5em}
西門送僕射，鞍馬照路光。\par
南門送貳卿，冠蓋遥相望。\par
東門送郎吏，艤舟隋堤傍。\par
郎吏誠隔品，同直白玉堂。\par
丈人况知己，振拔在舉場。\par
憐才惜我去，深勸離別觴。\par
醉中不記事，烟水空茫茫。\par
醒來聞鳴櫓，嘔軋摇斜陽。\par
猶疑在禁中，殘漏寒丁當。\par
迴望銀臺門，五雲遮帝鄉。\par
聚散本如此，升沉庸何傷。\par
丈人名位峻，只欠登巖廊。\par
平居倦朝請，高論思退藏。\par
圃田有別業，古木羅修篁。\par
艸亭寒蕭蕭，池波碧泱泱。\par
嘗云拂袖去，可以傲羲皇。\par
丈人果能爾，識度非尋常。\par
安車比疏廣，辟穀如張良。\par
再拜願丈人，壽考仍康强。\par
自念山野士，不解隨圓方。\par
宦途多齟齬，身計頗悲凉。\par
行將解簪笏，歸去事農桑。\par
幸容操杖履，洒掃近丘牆。\par

\section{送林介}
\textbf{王禹偁}\par\vspace{0.5em}
八年困名場，萬里有慈親。\par
吟變閩越聲，衣有京洛塵。\par
况復久相依，知子文行淳。\par
不能致一第，虛作金鑾臣。\par
昔予貶商洛，相送遠涉溱。\par
今予謫滁上，語別清淮濱。\par
途窮與道喪，詎免同沾巾。\par

\section{北樓感事}
\textbf{王禹偁}\par\vspace{0.5em}
北樓出林杪，登覽開病姿。\par
旁带滁州城，雉堞何逶迤。\par
下入刺史宅，却臨統軍池。\par
伊予翰林客，失職方在茲。\par
兩衙簿領外，盡日吟望時。\par
晚窗度急雨，夏木交繁枝。\par
淮南氣候殊，經秋囀黄鸝。\par
簷前有山果，採掇亦甘滋。\par
樽中有官醖，傾酌任醇醨。\par
忘機得真趣，懷古生遠思。\par
念昔李太尉，落落邦家基。\par
下筆到西漢，料兵如六奇。\par
謫官來此郡，鬰鬰持一麾。\par
嘗作懷嵩樓，記文悲盛衰。\par
甚得進退理，深明禍福機。\par
未幾再入用，斯言忽如遺。\par
君恩匪膠柱，天殃若影隨。\par
六月萬里行，炎荒竟不歸。\par
功成又名遂，不退將安之。\par
姑以人事較，勿憑天命推。\par
矧予草澤士，被褐復羹藜。\par
謬因弄文翰，八載侍丹墀。\par
三入承明盧，古人期並馳。\par
玉堂百日罷，所累非文詞。\par
强仕未爲老，望郎不爲卑。\par
淮邊永陽郡，人物自熙熙。\par
費用量所入，豐約從其宜。\par
一妻本糟糠，不識金翠施。\par
三男無庶孽，詎愛紈綺貲。\par
甘貧絕誅求，易退無覊縻。\par
五十擬歸耕，何必懸車期。\par
且予望衛公，雲龍與山麋。\par
唐賢昔際遇，文雅道光輝。\par
進士取將相，易于俯拾棋。\par
自從五代來，素風已陵遲。\par
干戈爲政事，茅土輸健兒。\par
儒冠筮仕者，僅免寒與饑。\par
至今明聖代，此風猶未移。\par
自無經濟術，烏用碌碌爲。\par
歸歟復歸歟，無忘北樓詩。\par

\section{送劉職方}
\textbf{王禹偁}\par\vspace{0.5em}
罷直出玉堂，詔授尚書郎。\par
朝僉假郡印，承乏來永陽。\par
喜開黄紙書，云替劉職方。\par
職方老狀頭，疇昔名擅場。\par
宦游三十載，寸進鴛鷺行。\par
未出六品班，鬚鬢已垂霜。\par
顧予新進生，遭時恣翱翔。\par
綸闈與內署，掌誥侍玉皇。\par
夜召赴浴殿，春遊醉柏梁。\par
晚出銀臺門，朝服含天香。\par
爲儒能到此，自道摩穹蒼。\par
今年四十二，典郡清淮旁。\par
卧錦郎位正，腰金服色光。\par
見君深自愧，撫己庸何傷。\par
古人重交代，子孫不相忘。\par
西掖替左司，劉白形詩章。\par
况我忝翰林，而君列文昌。\par
待將官與氏，山寺題修篁。\par
共保不朽名，他年比甘棠。\par
昇沉何足道，聚散安有常。\par
且當唱此詩，爲君送離觴。\par

\section{感興}
\textbf{王禹偁}\par\vspace{0.5em}
吾嘗入深山，谿谷寒且沍。\par
杉檜頗凌雲，歲月自朽蠹。\par
般輸目不見，何由用斤斧。\par
東山大夫松，中岳金雞樹。\par
秦政本獨夫，則天乃淫嫗。\par
名號被常材，所幸因一顧。\par
爲木豈有命，偶然生要路。\par
誰取澗底松，立作明堂柱。\par

\section{老態}
\textbf{王禹偁}\par\vspace{0.5em}
白髮不相饒，秋來生鬢邊。\par
黑花最相親，終日在眼前。\par
老態固具矣，宦情信悠然。\par
唯當共心約，收拾早歸田。\par

\section{官醖}
\textbf{王禹偁}\par\vspace{0.5em}
爲郡得官醖，月給盈三斛。\par
地僻少使車，時清罕留獄。\par
老大復遷謫，吾懷頗幽獨。\par
嬋娟樓上月，爛漫池邊菊。\par
東院與西亭，翛翛風弄竹。\par
對此不開樽，騷人應慟哭。\par
獨酌入醉鄉，陶然瞑雙目。\par
醒來成浩歎，胡爲事口腹。\par
彝酒書垂誡，羣飲聖所戮。\par
漢文亦禁酒，患在糜人穀。\par
自從孝武來，用度常不足。\par
榷酤奪人利，取錢入官屋。\par
古今事相倍，帝皇道難復。\par
吾無奈爾何，更盡杯中淥。\par

\section{射弩}
\textbf{王禹偁}\par\vspace{0.5em}
蹷張見舊史，强弩亦古官。\par
如何壯夫事，今作儒者歡。\par
罰郡在僻左，時清政多閑。\par
戎裝命賓侶，作此開愁顔。\par
吳弩號健捷，僕夫爲吾彎。\par
正侯廢已久，畫紙爲雕盤。\par
日暈生幾重，挂壁何團團。\par
記籌鼓聲急，中的酒量寬。\par
誠非軍旅事，亦有堵牆觀。\par
安得十萬枝，長驅過桑乾。\par
射彼老上庭，奪取燕脂山。\par
不見一匈奴，直抵瀚海還。\par
北方盡納款，獻壽天可汗。\par
吾徒久不武，干祿爲饑寒。\par
所得纔升斗，齷齪在朝班。\par
不如執戈士，意氣登韓壇。\par
笑擁白玉妓，醉馳黄金鞍。\par
郎官一生俸，供爾數月間。\par
儒將古所重，料兵如轉丸。\par
誰知戴章甫，終老弄辛酸。\par
偶因兒戲爲，痛念邊事艱。\par
臨風自慷慨，素髮衝儒冠。\par

\section{黑裘}
\textbf{王禹偁}\par\vspace{0.5em}
野蠶自成繭，繰絡爲山紬。\par
此物産何許，萊夷負海州。\par
一端重數斤，裁染爲吾裘。\par
守黑异華楚，示儉非輕柔。\par
燻香則無取，風雪曾何憂。\par
朝可奉冠带，夜以爲衾裯。\par
晏嬰三十年，庶幾跡相侔。\par
季子苦貂弊，吾服仍爲優。\par
不恥狐貉者，亦當師仲由。\par
况我屢遷謫，行將耕故丘。\par
映髮垂鷺頂，植杖昂鳩頭。\par
袖寬可以舞，老農即爲儔。\par
紫綬挂君門，任爾爭封侯。\par

\section{聞鴞}
\textbf{王禹偁}\par\vspace{0.5em}
元精育萬彙，羽族何茫茫。\par
爲怪有鴟鴞，爲瑞稱鳳皇。\par
鳳皇不時出，未識五色章。\par
吾生在鄒魯，風土殊遠方。\par
鳴鳩隨乳燕，日夕巢我梁。\par
翩翾雜鳥雀，穿屋率爲常。\par
又從筮仕來，五年居帝鄉。\par
更直入承明，侍宴趨未央。\par
上林聞鶑囀，巧舌如笙簧。\par
鴟鴞徒知名，聞見實未嘗。\par
頃年謫商山，聽之已悲凉。\par
今茲出內庭，罰郡來永陽。\par
誰知爾鸋鴂，營巢在城牆。\par
年加睡漸少，秋盡漏且長。\par
鳴嘯殊不已，歷歷含微霜。\par
孺人泣我右，稚子啼我傍。\par
吾心非達士，詎免亦倀倀。\par
人生縱百歲，忽若石火光。\par
其間有窮通，幽昧難自量。\par
我愛臯與夔，峨冠虞舜堂。\par
簫韶聞九成，丹穴來鏘鏘。\par
又愛閎與散，陳力遇文王。\par
鸑鷟聽岐山，多士周道昌。\par
嗟嗟漢賈誼，年少謫南荒。\par
故有鵩鳥賦，倚伏理甚詳。\par
郇公暨鄴侯，放逐同一邦。\par
夜深聞此鳥，韋公涕霑裳。\par
李侯舉酒令，斯音非不祥。\par
坐客如不聞，罰之以巨觴。\par
遂使惡鳥聲，聽之靡所傷。\par
贊皇貶袁州，懷鴞義亦臧。\par
乃知昔賢哲，未免亦悽遑。\par
况予不肖者，冒寵登朝行。\par
報國惟直道，謀身昧周防。\par
四年兩度黜，鬢髮已蒼蒼。\par
雖得五品官，銷盡百煉鋼。\par
何當解印綬，歸田謝膏粱。\par
教兒勤稼穡，與妻甘糟糠。\par
鳳來非我慶，鴞集非吾殃。\par
優游盡天年，身世俱可忘。\par

\section{甘菊冷淘}
\textbf{王禹偁}\par\vspace{0.5em}
經年厭粱肉，頗覺道氣渾。\par
孟春奉齋戒，敕厨唯素飱。\par
淮南地甚暖，甘菊生籬根。\par
長芽觸土膏，小葉弄晴暾。\par
采采忽盈把，洗去朝露痕。\par
俸麵新且細，溲攝如玉墩。\par
隨刀落銀鏤，煮投寒泉盆。\par
雜此青青色，芳草敵蘭蓀。\par
一舉無孑遺，空媿越盌存。\par
解衣露其腹，稚子爲我捫。\par
飽慚廣文鄭，饑謝魯山元。\par
况吾草澤士，藜藿供朝昏。\par
謬因事筆硯，名通金馬門。\par
官供政事食，久直紫微垣。\par
誰言謫滁上，吾族飽且溫。\par
既無甘旨慶，焉用品味繁。\par
子美重槐葉，直欲獻至尊。\par
起予有遺韵，甫也可與言。\par

\section{霪雨中偶書所見}
\textbf{王禹偁}\par\vspace{0.5em}
丙申二月七，是夕月離畢。\par
春雲忽霮䨴，春雨復蒙密。\par
綿綿殆三旬，不見天上日。\par
謫官在淮甸，幽抱常鬰鬰。\par
豈無瑯邪山，泥濘不可出。\par
空庭唯蚯蚓，得勢互蟠屈。\par
乘茲積陰氣，小穴恣出沒。\par
山禽忽飛下，長嘴啄深窟。\par
倒曳方力爭，强吞遽全失。\par
食土與巢林，上下非儔匹。\par
胡然羅此酷，不能保微質。\par
嗟嗟彼羣小，傾奪事匪一。\par
儀鳳去九夷，神龍入泉室。\par
如何役吾眼，瞻視此小物。\par
歸來因浩歎，自悔成詩筆。\par

\section{唱山歌}
\textbf{王禹偁}\par\vspace{0.5em}
滁民带楚俗，下俚同巴音。\par
歲稔又時安，春來恣歌吟。\par
接臂轉若環，聚首叢如林。\par
男女互相調，其詞非奔淫。\par
修教不易俗，吾亦弗之禁。\par
夜闌尚未闋，其樂何愔愔。\par
用此散楚兵，子房謀計深。\par
乃知國家事，成敗因人心。\par

\section{酬楊遂}
\textbf{王禹偁}\par\vspace{0.5em}
楊君江左士，文律何飄飄。\par
人言未冠時，作賦凌洞簫。\par
甲科中南國，通籍趨東朝。\par
轗軻位不進，陶潜還折腰。\par
宰邑向蜀道，龿蒲忽興妖。\par
官小力不支，奔竄避槍刀。\par
朝廷責守土，黜入縣佐僚。\par
昨朝寫孤憤，遺我有客謠。\par
伊予亦左遷，諷之心無憀。\par
人生一世間，否泰安可逃。\par
姑問道何如，未必論卑高。\par
自古富貴者，撩亂如藜蒿。\par
德業苟無取，未死名已消。\par
豈期顔子淵，不朽在一瓢。\par
推此任窮達，其樂方陶陶。\par
達則爲鵾鵬，窮則爲鷦鷯。\par
垂天與巢林，識分皆逍遥。\par

\section{和楊遂賀雨}
\textbf{王禹偁}\par\vspace{0.5em}
我罷內庭職，出臨永陽民。\par
永陽民雖庶，未免多饑貧。\par
富之既無術，齪齪爲謹身。\par
可堪今夏旱，如燎復如焚。\par
厥田本塗泥，坐見生埃氛。\par
稚老無所訴，嗷嗷望穹旻。\par
食祿憂人憂，蚤夜眉不伸。\par
促决獄中囚，徧禱境內神。\par
楚辭有山鬼，廟貌羅水濱。\par
胡法有浮圖，寺宇連城闉。\par
齋莊命寮寀，供給抽俸緡。\par
鼓笛迎湫水，香花照金輪。\par
誠知非典故，且慰旱熯人。\par
偶與天雨會，霶霶四郊勻。\par
插秧復修堰，野叟何欣欣。\par
可辦官賦調，亦免農艱辛。\par
燮調頼時相，感應由聖君。\par
于吾復何有，敢望歌頌云。\par
清流楊水部，德與我爲鄰。\par
仇香官位屈，何遜詩格新。\par
見投賀雨篇，言自人口聞。\par
夫君蓋私我，過實豈相親。\par
爲霖非我事，職業唯詞臣。\par
若有民謠起，當歌帝澤春。\par
庶使採詩官，入奏助南薰。\par

\section{揚州寒食贈屯田張員外成均吳博士同年殿省柳丞}
\textbf{王禹偁}\par\vspace{0.5em}
前年寒食節，待詔直內庭。\par
休假百官出，獨掩深嚴扃。\par
近侍不敢醉，賜酒空滿瓶。\par
閑就通中枕，時聞索上鈴。\par
思入無何鄉，兀然欲忘形。\par
去年寒食日，滁上忝專城。\par
山歌喧里巷，春物媚池亭。\par
永陽溪水淥，琅邪山色青。\par
謫宦自消遣，不敢誇獨醒。\par
往往取官醖，時時對花傾。\par
醉來念身世，翻使淚縱橫。\par
今年蒞淮海，時節又清明。\par
堆案有留事，聽歌無歡聲。\par
胥徒費簿領，使客煩送迎。\par
狴牢未空歇，堰埭勞修營。\par
衰病力不支，懶慢性已成。\par
虛花滿雙目，素髮添幾莖。\par
酒肴略無味，妓樂固難聽。\par
誰言寒食下，終日取茶烹。\par
屯田布素交，屈此關市征。\par
昔年同應舉，典衣飛巨觥。\par
博士東觀客，求官得步兵。\par
况且丹陛前，同爲出谷鶑。\par
殿丞尹我邑，桑梓復弟兄。\par
吏隠掌鹺茗，終朝看道經。\par
三賢宴會少，七夕休假并。\par
何不策我馬，廢苑尋放螢。\par
何不蕩我舟，樓基訪摘星。\par
三春景欲盡，九曲波始平。\par
居然逼吏役，頓此阻交情。\par
老態厭春華，病身憂宿酲。\par
如水若不改，藉糟亦胡寧。\par
解印蓄素志，吟詩露丹誠。\par
維揚非所愛，有便即歸耕。\par

\section{揚州池亭即事}
\textbf{王禹偁}\par\vspace{0.5em}
冥心閱羣動，亦各趨所安。\par
胡爲名利人，戚戚常鮮歡。\par
吾生四十四，結佩呼郎官。\par
掌言入綸閣，待詔直金鑾。\par
匪謂得祿少，所嗟行道難。\par
前年謫滁上，憂畏雙鬢殘。\par
頼有琅邪溪，時濯塵纓冠。\par
朝僉徙淮海，任重力易殫。\par
君恩詎可報，感激涕汍瀾。\par
民瘼不能治，惻隠情悲酸。\par
況復多病身，名宦心已闌。\par
歸田未果决，懷祿尚盤桓。\par
公退何所適，池亭一凭欄。\par
旭日媚春卉，微風生鳴湍。\par
呵僮勿挾彈，留客不持竿。\par
用冀魚鳥馴，熙熙肆遊觀。\par
神仙未可學，吏隠聊自寬。\par
孤吟刻幽石，此義非考槃。\par

\section{一品孫鄭昱}
\textbf{王禹偁}\par\vspace{0.5em}
卜葬得假告，南出安上門。\par
鞭馬六十里，暮投中書村。\par
村翁館我宿，茅屋欲黄昏。\par
有客忽投刺，自稱一品孫。\par
氣貌不凡俗，因爲開酒樽。\par
坐久問家諜，其族大且繁。\par
池州有清節，濫觴登洪源。\par
太傅擅鴻筆，入相又出藩。\par
其家本開封，改號一何尊。\par
至昱始六代，布衣老丘樊。\par
跨馿入府縣，驅犢耕郊原。\par
家廟固已毀，國史空具存。\par
盛德百世著，功必格乾坤。\par
高太已不祀，羨絪何可論。\par
况復起章句，乘時寵便蕃。\par
子孫雖替陵，尚得守田園。\par
我愛三代時，法度有深根。\par
卿大夫稱家，世世奉蘋蘩。\par
四民有定分，宦路無馳奔。\par
自從雜伯道，傾奪日喧喧。\par
脫耒秉金鉞，吮筆乘朱軒。\par
朝榮又暮辱，容易如掌翻。\par
古道不可復，頹波益以渾。\par
何况度木者，倒置輪與轅。\par
我亦起白屋，兩朝直紫垣。\par
蔭子有官常，賞延弟與昆。\par
盡待食人祿，將何報君恩。\par
農桑國之本，孝義古所敦。\par
吾族不力穡，終歲飽且溫。\par
雖非享富貴，亦以蠹黎元。\par
唐賢尚消歇，我輩奚足言。\par
呼兒諷此詩，播在篪與塤。\par

\section{鳳皇陂}
\textbf{王禹偁}\par\vspace{0.5em}
次公治潁川，仁政被一方。\par
神物不藏瑞，茲焉集鳳皇。\par
在昔奏簫韶，舜庭來蹌蹌。\par
西伯有至化，亦見鳴岐陽。\par
仲尼豈無德，已矣空悲傷。\par
夫何刀筆吏，而能致殊祥。\par
我來過荒陂，烟草但蒼蒼。\par
緬懷漢循吏，史筆恐未詳。\par

\section{月波樓詠懷}
\textbf{王禹偁}\par\vspace{0.5em}
郡城無大小，雉堞皆有樓。\par
其間著名者，不過十數州。\par
吹簫事遼夐，仙跡難尋求。\par
庾公在九江，締構何風流。\par
謝守鎮宣城，疊嶂名有由。\par
東陽敞八詠，吾聞沈隠侯。\par
白雪架郢中，調高難和酬。\par
黄鶴倚鄂渚，仙去事悠悠。\par
贊皇謫滁上，作賦懷嵩丘。\par
樓居出俗態，澤國多勝遊。\par
好景不遇人，安得名存留。\par
齊安古郡廢，移此清江頭。\par
築城隨山勢，屈曲復環周。\par
茲樓最軒豁，曠望西北陬。\par
武昌地如掌，天末入雙眸。\par
平遠無林木，一望同離婁。\par
山形如八字，會合勢相勾。\par
三國事既遠，六朝名亦休。\par
近從唐末來，爭奪互仇讎。\par
斯樓備矢石，此地控咽喉。\par
終朝望烽燧，連歲事戈矛。\par
可憐好詩景，牢落無人收。\par
皇家統萬國，遠邇盡懷柔。\par
三聖四十年，蕩蕩文德修。\par
淮甸爲內地，黄岡壓上游。\par
儒冠假郡印，踐更若公郵。\par
況多辦職吏，誰肯恣吟謳。\par
伊余何爲者，竊慕騷人儔。\par
兩朝掌文翰，十年侍冕旒。\par
去歲出西掖，謫居抱窮愁。\par
日日江樓上，風物得冥搜。\par
何人名月波，此義頗爲優。\par
西南新桂魄，初上懸玉鈎。\par
曉瀨清且淺，漂蕩影沉浮。\par
三五金波滿，夜光如暗投。\par
驪龍弄頷珠，晃朗照汀洲。\par
澹臺拔寶劍，碎璧斬長虬。\par
冰輪曉入地，推下赤金毬。\par
闌干四五星，斜漢印清秋。\par
誰家上元燈，兒戲刳𦸈𧁾。\par
此景吟不出，謾使聲呦呦。\par
千里畫圖闊，四時詩興幽。\par
野花媚宮纈，芳草鋪碧紬。\par
火雲照沙浦，暴雨傾瓦溝。\par
白亂蘆花散，紅殷蓼穗稠。\par
簷冰垂若綆，雪片大于鷗。\par
江蘺煙漠漠，官柳雨颼颼。\par
舟子斜盪槳，牧童倒騎牛。\par
水獺有時戲，江豚頗能泅。\par
山鳥奏竽籟，落霞展衾裯。\par
魚網雪離離，酒旗風飂飂。\par
旅懷雖自適，詩物奈相尤。\par
右顧徐邈洞，精靈知在否。\par
左瞰伍員廟，荒隙令人羞。\par
樓中何所有，官醖湛蚍蜉。\par
棋枰留客坐，琴調待僧抽。\par
橘苞鄰藥鼎，詩筆間茶甌。\par
平生性幽獨，寂寞誰獻酬。\par
官常已三黜，懷抱罹百憂。\par
凭欄憶王粲，望闕同子牟。\par
自甘成潦倒，無復事聲猷。\par
身世喻泡幻，衣冠如贅瘤。\par
放意無何鄉，誰分親與仇。\par
寓形朝籍中，毀譽任啁啾。\par
君恩無路報，民瘼無術瘳。\par
唯慚戀祿俸，未去耕田疇。\par
題詩郡樓上，含毫思夷猶。\par
功名非范蠡，何必泛扁舟。\par

\section{十月二十日作}
\textbf{王禹偁}\par\vspace{0.5em}
重衾又重茵，蓋覆衰孏身。\par
中夜忽涕泗，無復及吾親。\par
須臾殘漏歇，吏報國忌辰。\par
凌旦騎馬出，溪冰薄潾潾。\par
路傍饑凍者，顔色頗悲辛。\par
飽暖我不覺，羞見黄州民。\par
昔賢終祿養，往往歸隠淪。\par
誰教爲妻子，頭白走風塵。\par
修身與行道，多愧古時人。\par

\section{雷}
\textbf{王禹偁}\par\vspace{0.5em}
商山春夏旱，旱雷不降雨。\par
及秋又霖霪，雷聲時一舉。\par
南山復北山，渹礚若贙虎。\par
君子容必變，所以敬天怒。\par
無乃豐隆鬼，恣意撾雷鼓。\par
此怒既非天，敬之亦奚取。\par

\section{秋霖二首  其一}
\textbf{王禹偁}\par\vspace{0.5em}
秋霖過百日，歲望終何如。\par
嘉穀就穗生，茁茁垂青鬚。\par
宿麥未入土，大田多泥塗。\par
河闊不辨馬，原高恐生魚。\par
時政苟云失，生民亦何辜。\par
雨若是天淚，天眼應已枯。\par

\section{秋霖二首  其二}
\textbf{王禹偁}\par\vspace{0.5em}
山雲百日雨，山水十丈波。\par
田疇與道路，一夕成江河。\par
巨石大於甕，吹轉如蓬窠。\par
夏旱既損麥，秋潦復無禾。\par
津梁盡傾壞，商販絕經過。\par
斗米二百金，吾生將奈何。\par
安敢比夷齊，愚聖不同科。\par
應如元魯山，餓死深山阿。\par

\section{真娘墓}
\textbf{王禹偁}\par\vspace{0.5em}
女命在于色，士命在于才。\par
無色無才者，未死如塵灰。\par
虎丘真娘墓，止是空土堆。\par
香魂與膩骨，銷散隨黄埃。\par
何事千百年，一名長在哉。\par
吳越多婦人，死即藏山隈。\par
無色固無名，丘冢空崔嵬。\par
唯此真娘墓，客到情徘徊。\par
我是好名士，爲爾傾一杯。\par
我非好色者，後人無相咍。\par

\section{遊虎丘}
\textbf{王禹偁}\par\vspace{0.5em}
樂天曾守郡，酷愛虎丘山。\par
一年十二度，五馬來松關。\par
我今方吏隠，心在雲水間。\par
野性羣麋鹿，忘機狎鷗鷳。\par
乘興即一到，興盡復自還。\par
不知使君貴，何似長官閑。\par

\section{吳王墓}
\textbf{王禹偁}\par\vspace{0.5em}
惜哉吳王墓，秦帝嘗開破。\par
應笑埋金玉，千年賈餘禍。\par
不待虎迹銷，已聞鮑車過。\par
又是驪山頭，炎炎三月火。\par

\section{橄欖}
\textbf{王禹偁}\par\vspace{0.5em}
江東多果實，橄欖稱珍奇。\par
北人將就酒，食之先顰眉。\par
皮核苦且澀，歷口復棄遺。\par
良久有迴味，始覺甘如飴。\par
我今何所喻，喻彼忠臣詞。\par
直道逆君耳，斥逐投天涯。\par
世亂思其言，噬臍焉能追。\par
寄語採詩者，無輕橄欖詩。\par

\section{爲惡}
\textbf{王禹偁}\par\vspace{0.5em}
明時巧言士，亂世佞倖郎。\par
佞倖惟苟且，巧言頗包藏。\par
爲惡雖不同，同歸於覆亡。\par

\section{成武縣作}
\textbf{王禹偁}\par\vspace{0.5em}
釋褐來成武，初官且自强。\par
位卑松在澗，俸薄葉經霜。\par
雨菌生書案，饑禽啄印床。\par
猶驚寫秋卷，槐砌落花黄。\par

\section{寄碭山主簿朱九齡}
\textbf{王禹偁}\par\vspace{0.5em}
閑思蓬島會羣仙，二百同年最少年。\par
利市襴衫拋白紵，風流名紙寫紅牋。\par
歌樓夜宴停銀燭，柳巷春泥污錦韉。\par
今日折腰塵土裏，共君追想好淒然。\par

\section{寄魚臺主簿傅翶}
\textbf{王禹偁}\par\vspace{0.5em}
健羨魚臺景最奇，鮑參軍到語多時。\par
波平綠野懸魚網，木脫空城露酒旗。\par
錦擲鮮鱗紅撥剌，雪翹寒鷺白䙰褷。\par
仍誇縣尹風騷客，應有秋來唱和詩。\par

\section{寄寧陵陳長官}
\textbf{王禹偁}\par\vspace{0.5em}
吏隠寧陵縣，琴堂枕古河。\par
家山隔江遠，風雨過船多。\par
假日親尋藥，公庭自種莎。\par
相逢如舊識，執手動勞歌。\par

\section{寄金鄉張贊善}
\textbf{王禹偁}\par\vspace{0.5em}
年少辭榮自古稀，朝衣不著著斑衣。\par
北堂侍膳侵星起，南畝催耕冒雨歸。\par
種竹野塘春筍脆，採蘭幽澗露牙肥。\par
伊予自是徒勞者，未得同尋舊釣磯。\par

\section{赴長洲縣作  其一}
\textbf{王禹偁}\par\vspace{0.5em}
移任長洲縣，扁舟興有餘。\par
篷高時見月，棹穩不妨書。\par
雨碧蘆枝亞，霜紅蓼穗疏。\par
此行紆墨綬，不是爲鱸魚。\par

\section{赴長洲縣作  其二}
\textbf{王禹偁}\par\vspace{0.5em}
移任長洲縣，辭親淚滿衣。\par
折腰雖未免，搔首欲何歸。\par
曉白霜華重，晴紅栗葉飛。\par
江頭鷗鳥在，應怪不忘機。\par

\section{赴長洲縣作  其三}
\textbf{王禹偁}\par\vspace{0.5em}
移任長洲縣，孤帆冒雨行。\par
全家隨逆旅，一夜泊江城。\par
身世漂淪極，功名早晚成。\par
惟當泥尊酒，得喪任浮生。\par

\section{赴長洲縣作  其四}
\textbf{王禹偁}\par\vspace{0.5em}
移任長洲縣，窮秋入水鄉。\par
江涵千頃月，船載一篷霜。\par
竹密藏魚市，雲疏漏雁行。\par
故園漸迢遞，烟浪白茫茫。\par

\section{赴長洲縣作  其五}
\textbf{王禹偁}\par\vspace{0.5em}
移任長洲縣，沿流漸入吳。\par
見碑時下岸，逢店自徵酤。\par
野廟連荒冢，江禽似畫圖。\par
高堂從別後，應夢宿菰蒲。\par

\section{惠山寺留題}
\textbf{王禹偁}\par\vspace{0.5em}
吟入惠山山下寺，古泉閑挹味何嘉。\par
好拋此日陶潛米，學煮當年陸羽茶。\par
猶負片心眠水石，略開塵眼識烟霞。\par
勞生未了還東去，孤棹寒篷宿浪花。\par

\section{遊虎丘寺}
\textbf{王禹偁}\par\vspace{0.5em}
寺牆圍著碧孱顔，曾是當年海湧山。\par
盡抱好峰藏院裏，不教幽景落人間。\par
劍池草色經冬在，石座苔花自古斑。\par
珍重晉朝吾祖宅，一迴來此便忘還。\par

\section{寄獻潤州趙舍人  其一}
\textbf{王禹偁}\par\vspace{0.5em}
南徐城古樹蒼蒼，衙府樓臺盡枕江。\par
甘露鐘聲清醉榻，海門山色滴吟窗。\par
直廬久負題紅葉，出鎮何妨擁碧幢。\par
聞說秋來自高尚，道裝筇竹鶴成雙。\par

\section{寄獻潤州趙舍人  其二}
\textbf{王禹偁}\par\vspace{0.5em}
記言彩筆罷摛華，郡閣高閑似道家。\par
琴院坐聽江寺磬，郡樓吟見海山霞。\par
春園遺母親燒筍，夜榻留僧自煮茶。\par
應笑陶潛未歸去，折腰奔走在泥沙。\par

\section{寄毗陵劉博士}
\textbf{王禹偁}\par\vspace{0.5em}
毗陵古郡接江堧，赴任琴書共一船。\par
下岸且尋甘露寺，到城先問惠山泉。\par
秋蟾吐檻供吟興，野鶴偎床伴醉眠。\par
官散道孤詩格老，算應雙鬢更皤然。\par

\section{謝柴侍御送鶴}
\textbf{王禹偁}\par\vspace{0.5em}
兩朵春雲墮世間，霜臺乞與縣僚看。\par
骨含仙氣生來瘦，羽插天風過處寒。\par
唳別繡衣聲寂寞，舞隨藍綬意闌珊。\par
今朝窗外頻相顧，應怪頭無獬豸冠。\par

\section{官舍書懷呈羅思純}
\textbf{王禹偁}\par\vspace{0.5em}
同年事分幾般同，墨綬逶迤一郡中。\par
仙桂並枝攀月魄，縣花交影笑春風。\par
分莎種就尋僧逕，借竹編成養鶴籠。\par
公暇不妨閑唱和，免教來往遞詩筒。\par

\section{除夜寄羅評事同年  其一}
\textbf{王禹偁}\par\vspace{0.5em}
歲暮洞庭山，知君思浩然。\par
年侵曉色盡，人枕夜濤眠。\par
移棹燈摇浪，開窗雪滿天。\par
無因一乘興，同醉太湖船。\par

\section{除夜寄羅評事同年  其二}
\textbf{王禹偁}\par\vspace{0.5em}
除夜在天涯，共君同憶歸。\par
夢中傾壽酒，堂下著斑衣。\par
臘雪通宵舞，梅花到處飛。\par
更堪洞庭宿，愁思倍依依。\par

\section{除夜寄羅評事同年  其三}
\textbf{王禹偁}\par\vspace{0.5em}
郡僚方賀正，獨宿太湖稜。\par
階下羞爲吏，船中祗載僧。\par
折梅和薄雪，煮茗對孤燈。\par
應笑排衙早，寒靴踏曉冰。\par

\section{陸羽泉茶}
\textbf{王禹偁}\par\vspace{0.5em}
甃石封苔百尺深，試茶嘗味少知音。\par
唯餘半夜泉中月，留得先生一片心。\par

\section{投柴殿院}
\textbf{王禹偁}\par\vspace{0.5em}
南面修文德，東吳納土疆。\par
蒼生思撫育，丹詔擇循良。\par
烏府官新轉，龍頭桂舊香。\par
渡江驄馬瘦，垂地繡衣長。\par
綸閣材知屈，蘇臺俗必康。\par
恩流一車雨，威凜柏臺霜。\par
休假尋山寺，行春泊野塘。\par
白公是前政，魯望有維桑。\par
求瘼心雖切，頤神道豈妨。\par
煎茶虎丘井，搗藥木蘭堂。\par
筍蕨供家饌，園林著道裝。\par
擊筇教鶴舞，收橘待僧嘗。\par
迎使朝衣穩，娛賓綺席張。\par
犬聲銷巷陌，鶑舌動笙簧。\par
冷句題秋葉，孤琴貯夜囊。\par
歌樓寒月白，飲舫晚波凉。\par
官業除苛法，家風襲雅章。\par
豸冠危肅物，象簡醉橫床。\par
熊軾淹寧久，鼇頭譽轉芳。\par
南園休命侶，北闕即徵黄。\par
清貴容誰見，遭逢合自强。\par
字人叨屬邑，畏德每循牆。\par
名品知懸隔，孤危俟薦揚。\par
折腰休太息，青眼異尋常。\par
岱岳容拳石，滄溟納濫觴。\par
扶摇如借便，羽翼必高翔。\par
從事員多闕，徒勞跡可傷。\par
金臺雖載築，珠履未成行。\par
始隗前言在，依劉後進光。\par
免教青史上，徒美一燕王。\par

\section{東鄰竹}
\textbf{王禹偁}\par\vspace{0.5em}
東鄰誰種竹，偏稱長官心。\par
月上分清影，風來惠好音。\par
低枝疑見接，迸筍似相尋。\par
多謝此君意，牆頭誘我吟。\par

\section{南園偶題}
\textbf{王禹偁}\par\vspace{0.5em}
天子優賢是有唐，鑑湖恩賜賀知章。\par
他年我若功成去，乞取南園作醉鄉。\par

\section{戲贈嘉興朱宰同年}
\textbf{王禹偁}\par\vspace{0.5em}
縣前蘇小舊荒墳，應作行雲入夢頻。\par
猶勝白公尉盩厔，庭花剛喚作夫人。\par

\section{春日官舍偶題}
\textbf{王禹偁}\par\vspace{0.5em}
薄宦苦流離，壯年心已衰。\par
鶑花愁不覺，風雨病先知。\par
曉月晃竹屋，寒苔疊槿籬。\par
無人慰幽寂，庭柳自低垂。\par

\section{寄獻翰林宋舍人}
\textbf{王禹偁}\par\vspace{0.5em}
金鼎鹽梅偶未和，位高猶說野情多。\par
宮牆月上開琴匣，道院風清響藥羅。\par
留客旋燒含露筍，倩僧教種耐霜莎。\par
孤寒知有爲霖望，未忍江頭釣淥波。\par

\section{吳江縣寺留題}
\textbf{王禹偁}\par\vspace{0.5em}
松江江寺對峰巒，檻外生池接野灘。\par
幽鷺靜翹春草碧，病僧閑說夜濤寒。\par
晨齋施筍唯溪叟，國忌行香祇縣官。\par
盡日門前照流水，塵纓渾疑濯汍瀾。\par

\section{獻轉運使雷諫議  其一}
\textbf{王禹偁}\par\vspace{0.5em}
當年辭氣壓朱雲，老作皇家諫諍臣。\par
章疏罷封無事日，朝廷猶惜直言人。\par
題詩野館光泉石，講易秋堂動鬼神。\par
棘寺下僚叨末路，齋心唯願秉陶鈞。\par

\section{獻轉運使雷諫議  其二}
\textbf{王禹偁}\par\vspace{0.5em}
江南江北接王畿，漕運帆檣去似飛。\par
父子有才同富國，君王無事免宵衣。\par
屏除姦吏魂應喪，養活疲民肉漸肥。\par
還有文場受恩客，望塵情抱倍依依。\par

\end{document}